%% P04 --- Vectors, vector spaces and basis
%%
%% Every number in this program that the reader cannot do in their head comes
%% from code/p04_vector_spaces.py and is pulled in with \val{}. The console
%% listing is a committed transcript written by that script.
%%
%% THE THESIS: span, linear independence, basis and dimension are FOUR WORDS
%% FOR ONE IDEA, and dimension is the number of independent directions rather
%% than the length of the list.
%%
%% THE FIRST PROGRAM IN THIS BOOK NEEDING A NEW OBJECT. Everything before it
%% asked a new question about objects the reader already had; Program P03's
%% closing frame says so. Part III opens here.
%%
%% WHAT THIS PROGRAM IS OWED AND PAYS. Program F09 defined vector arithmetic
%% and then said, in a warning box, that how vectors are ARRANGED in high
%% dimension is a different class of question, that P05 measures it, and that
%% measuring it properly needs the vocabulary THIS program sets up. That
%% vocabulary is the whole deliverable.
%%
%% WHAT THIS PROGRAM DELIBERATELY LEAVES ALONE, checked against
%% tools/programs.json rather than remembered:
%%   inner products, norms, projection, and near-orthogonality in high
%%     dimension -- the measurement F09 deferred                    -> P05
%%   matrices as linear maps, multiplication as composition         -> P06
%%   rank as a property of a matrix, and LoRA                       -> P08
%% P04 never measures an angle, because it has not been given one. Section 5
%% goes as far as a COUNTING THEOREM and stops, and says in as many words that
%% the interpretability stories built on top of it are hypotheses.
%%
%% Twin: programs/pl/P04-vector-spaces.tex. Frame for frame, macro for macro,
%% number for number; tools/parity.py compares them.

\program{Vectors, vector spaces and basis}
\label{prog:P04}
\index{vector!space}
\index{basis}

\begin{outcomes}
\outcome{1--8}{say what makes a set of things a vector space, and test a given
  set against it}
\outcome{9--14}{say what a set of vectors spans, and what the span is not}
\outcome{15--20}{decide whether a list of vectors is independent, and say what
  bounds how many can be}
\outcome{21--27}{say what a basis is, what a dimension counts, and which of the
  two a set of coordinates belongs to}
\outcome{28--34}{apply the counting bound to an embedding table, and separate
  what it proves from what it does not}
\end{outcomes}

\begin{quiz}
\item What two operations does a vector space have, and what must be true of
  their results? \teachesat{1--2}
  \answerto{Addition of two vectors and multiplication by a scalar, and both
  results must land back inside the set. That closure is the whole of the
  definition that matters here.}
\item Is the set of vectors in $\R^{3}$ whose first component is $1$ a vector
  space? \teachesat{2--3}
  \answerto{No. Add two of them and the first component is $2$, which is
  outside the set. The set whose first component is $0$ is one.}
\item Two vectors in $\R^{3}$ are given. What do they span? \teachesat{10--11}
  \answerto{A plane through the origin \dash{} unless one is a multiple of the
  other, in which case a line. That exception is linear dependence, before it
  has been named.}
\item Can $\val{p04.family.count}$ vectors of $\val{p04.family.ambient}$
  components each be linearly independent? \teachesat{16--17}
  \answerto{No. At most $\val{p04.family.ambient}$ vectors can be independent
  in $\val{p04.family.ambient}$ dimensions, whatever they are and however many
  you write down.}
\item A vector's coordinates change when the basis changes. What does not
  change? \teachesat{24--25}
  \answerto{The vector. Its length, and its relations to every other vector,
  are properties of the space rather than of the axes you wrote it against.}
\item A model embeds a vocabulary of $\val{p04.vocab}$ tokens in
  $\val{p04.embed}$ dimensions. How many of the embeddings must be combinations
  of the others? \teachesat{28--29}
  \answerto{At least $\val{p04.dependent}$ of them, which is
  $\val{p04.dependent.pct}\%$. It is a counting theorem and no property of the
  training changes it.}
\item Does that mean the model cannot represent $\val{p04.vocab}$ distinct
  concepts? \teachesat{30--31}
  \answerto{No. It means it cannot give each of them its own independent
  direction. What it does instead is a question about structure, and the
  answers on offer are hypotheses.}
\end{quiz}

% ============================================================
\section{What a space is}
\label{sec:P04-space}
\index{vector!space}

\begin{fr}
Part II asked three questions about one arithmetic and needed no new objects to
do it. This program needs one, and it is the first in this book that does.

Program~\ref{prog:F09} gave you a vector and its arithmetic: a list of numbers,
added componentwise, scaled by a single number. It also left something
unfinished. Its closing warning says that how vectors are \emph{arranged} in
high dimension is a different class of question from how they are computed with,
that Program~\ref{prog:P05} measures it, and that measuring it needs vocabulary
this program has to supply first.

So start where Program~\ref{prog:F09} stopped. Given two vectors, what are you
allowed to do to them?
\nextframe
\end{fr}

\begin{fr}
\begin{ansblock}
Add them, and multiply one of them by a number.
\end{ansblock}

That is the whole list, and its shortness is the point. There is no
multiplication of two vectors here, no division, no ordering. Two operations.

A \emph{vector space} is a set on which those two operations are defined and
which is \textbf{closed} under both: add any two members and the answer is a
member; scale any member and the answer is a member. There are other axioms and
they are the ones you would write down without being asked \dash{} addition
commutes, scaling distributes \dash{} and closure is the one that ever fails in
practice.

Test it. Is the set of vectors in $\R^{3}$ whose first component is $0$ a
vector space? Is the set whose first component is $1$?
\nextframe
\end{fr}

\begin{fr}
\begin{ansblock}
The first is. The second is not.
\end{ansblock}

Add $(0, a, b)$ and $(0, c, d)$ and the first component is still $0$; scale one
and it is still $0$. Closed on both counts.

Add $(1, a, b)$ and $(1, c, d)$ and the first component is $2$, which is not in
the set. One counterexample and it is not a space \dash{} and notice that the
set is a perfectly reasonable set of vectors. It is a plane and it is
$\R^{3}$-shaped.

There is a quicker test than either closure check, and this set fails that one
too. Both properties above are about what stays inside a set; name the single
vector every space is therefore obliged to contain, and say which of the two
properties forces it.
\dotline
\nextframe
\end{fr}

\begin{fr}
\ans{The zero vector, forced by scaling: take any member and scale it by $0$.}

\textbf{Every vector space contains the zero vector}, so a set that does not is
not one, and the check costs nothing.

\begin{aibox}
The set of embeddings your model actually produces is not a vector space, and
the distinction is worth holding.

Twenty thousand rows of a table are twenty thousand vectors. The space they
live in is $\R^{\val{p04.embed}}$, and it contains every combination of them
whether the model ever emits one or not. Most of the interesting claims in
interpretability are about \emph{where in that space} the model's outputs lie
\dash{} which is a claim about a subset, and a subset with structure is what
those claims are trying to describe.

Confusing the two is how \enquote{the model has a direction for sentiment}
becomes ambiguous. A direction exists in the space by construction. Whether the
model \emph{uses} it is the question, and it is not settled by the space.
\end{aibox}
\end{fr}

\begin{fr}
One more thing about the definition, because it is what makes the vocabulary
worth learning rather than a formality.

Nothing in the two operations mentions lists of numbers. Polynomials can be
added and scaled, so they form a vector space. So do functions, so do infinite
sequences, so do matrices when you meet them.

\textbf{\enquote{Vector} is a role rather than a shape.} Everything the rest of
this program proves is proved from closure alone, so all of it applies to every
one of those, and none of it has to be re-derived when the objects change.
\end{fr}

\begin{fr}
A subset that is itself a space, under the same two operations, is a
\emph{subspace}. The first-component-zero set above is one; the
first-component-one set is not.

Subspaces are where the rest of this program lives, because the interesting
question is almost never \enquote{is this a space} \dash{} it is
\textbf{which part of the space is being used}.

Test one more, and this one is worth getting wrong. Is the set of vectors in
$\R^{3}$ with at most one non-zero component a subspace?
\nextframe
\end{fr}

\begin{fr}
\ans{No.}
\begin{trapbox}
\textbf{Closure has to hold for \emph{every} pair, and a set can be closed
under scaling and fail under addition.}

Scale $(a, 0, 0)$ by anything and it still has one non-zero component, so that
half is fine. But $(1, 0, 0) + (0, 1, 0) = (1, 1, 0)$, which has two, and the
set does not contain it.

This is the shape of nearly every failed closure argument: the scaling is
obvious, the addition is where the counterexample lives, and the set that fails
is usually one somebody described as \enquote{sparse}. The union of the three
axes is not a subspace; the plane containing two of them is.
\end{trapbox}
\end{fr}

\begin{fr}
The vocabulary so far is three words \dash{} space, subspace, closure \dash{}
and one test you can apply in your head.

What is missing is any way of saying how \emph{big} a subspace is. The three
axes and the plane through two of them are both subsets of $\R^{3}$, and calling
one bigger than the other needs a notion of size that counting the vectors
cannot supply: both contain infinitely many.

The next three sections build that notion, and they build it out of the two
operations and nothing else.
\end{fr}

% ============================================================
\section{Span}
\label{sec:P04-span}
\index{vector!span}

\begin{fr}
Take a handful of vectors and apply the two operations to them in every
possible way: scale each one by any number, add the results, repeat. The set of
everything you can reach is the \emph{span} of the original handful.

\medskip
\noindent
$\spanof\{v_1, \dots, v_k\}
  = \{\, a_1 v_1 + \dots + a_k v_k \;:\; a_i \text{ any numbers} \,\}$

\medskip
\noindent
A sum like that is a \emph{linear combination}, and it is the only construction
in this whole subject. Every definition below is a sentence about linear
combinations.

The span is always a subspace, and the reason is one line: add two linear
combinations of the $v_i$ and you get a linear combination of the $v_i$.
\end{fr}

\mermaidfig{p04-span}{The two operations, applied in every combination, and
what they reach.}{what a span is}

\begin{fr}
Two vectors in $\R^{3}$, neither of them zero. What do they span?
\nextframe
\end{fr}

\begin{fr}
\ans{A plane through the origin \dash{} usually.}
Through the origin, because taking both coefficients to be $0$ is one of the
combinations and it gives the zero vector. Every span contains it.

And \emph{usually}, which is the word to stop at. Two vectors span a plane in
almost every case, and there is one arrangement of them for which they do not.
Name it, and say what the span is instead.
\dotline
\nextframe
\end{fr}

\begin{fr}
\ans{When one is a multiple of the other. Every combination is then a multiple
of either, so the span is a line.}

That exception is the whole of the next section, arriving before it has a name.
\textbf{Two vectors do not always give you two directions.}

Which raises the question this program exists to answer. Here is a family of
$\val{p04.family.count}$ vectors, each with $\val{p04.family.ambient}$
components, and every one of them is built from the same two:
\[ v = a\,(1, 0, 2, -1, 3) + b\,(0, 1, -1, 4, 2) \]
for eight different pairs $(a, b)$.

How many dimensions does their span have?
\nextframe
\end{fr}

\begin{fr}
\ans{$\val{p04.family.dim}$.}
Not $\val{p04.family.count}$, which is how many vectors there are, and not
$\val{p04.family.ambient}$, which is how many components each one has. Every
member of the family is a combination of two fixed vectors, so every
combination of members is also a combination of those two, and the span is the
plane they determine.

Three numbers are now on the table and they are all different: how many
vectors are in the list, how long each vector is, and how many independent
directions they actually reach between them. Say which of the three is the
dimension.
\dotline
\nextframe
\end{fr}

\begin{fr}
\ans{The third. The first two are properties of how the data was written down;
only the third is a property of the data.}

That distinction is worth one more frame, because the words for the first two
get used for the third constantly.

An embedding table with $\val{p04.vocab}$ rows and $\val{p04.embed}$ columns
has three numbers attached to it in exactly the same way. Somebody will call it
\enquote{$\val{p04.embed}$-dimensional}, and they mean the second one \dash{}
the length of each row.

Whether the rows reach $\val{p04.embed}$ independent directions is a separate
question with a separate answer, and section 5 shows that the answer to a
related question is forced. What is needed first is the word for the property
itself.
\end{fr}

% ============================================================
\section{Independence}
\label{sec:P04-independence}
\index{vector!linear independence}

\begin{fr}
A list of vectors is \emph{linearly dependent} when at least one of them is a
linear combination of the others \dash{} when it adds nothing to the span, so
that dropping it changes nothing.

It is \emph{linearly independent} when no such vector exists.

There is a second form of the definition that looks stranger and is easier to
use: the vectors are independent when the only way to combine them into the
zero vector is to take every coefficient to be zero.

The two say the same thing, and it is worth seeing why. If some combination
$a_1 v_1 + \dots + a_k v_k = 0$ has a non-zero coefficient \dash{} say $a_1$
\dash{} then dividing by it makes $v_1$ a combination of the rest.
\end{fr}

\begin{fr}
Now the question the last section set up.

Can $\val{p04.family.count}$ vectors, each with $\val{p04.family.ambient}$
components, be linearly independent? Not \emph{are these eight}, but can any
eight such vectors be?
\nextframe
\end{fr}

\begin{fr}
\ans{No, and it is a theorem rather than a fact about those eight.}
\transcript{p04-rank}

The transcript takes random vectors in $\val{p04.sat.dim}$ dimensions and
counts how many independent directions each set reaches. Two vectors reach
$\val{p04.sat.rank.2}$, four reach $\val{p04.sat.rank.4}$, eight reach
$\val{p04.sat.rank.8}$ \dash{} and then twelve reach
$\val{p04.sat.rank.12}$, twenty reach $\val{p04.sat.rank.20}$, and forty reach
$\val{p04.sat.rank.40}$.

\textbf{It saturates, and it saturates at the number of components.}
\end{fr}

\mermaidfig{p04-saturation}{Why the count of genuinely new directions has to
stop rising.}{why the count saturates}

\begin{fr}
\begin{trapbox}
\textbf{More vectors do not give you more directions once you have run out of
directions to give.}

The reasoning behind the saturation needs no machinery. Each new vector either
points somewhere the earlier ones could not reach, in which case the span grows
by one dimension, or it is already a combination of them, in which case the
span does not grow at all. The span sits inside a space of
$\val{p04.sat.dim}$ dimensions, so it can grow at most
$\val{p04.sat.dim}$ times.

The script checks it over $\val{p04.sat.trials}$ random draws and never finds an
exception, which is the right kind of reassurance and the wrong kind of
evidence.
\end{trapbox}

Say what those draws are evidence of, given that the statement they are
checking is a theorem.
\dotline
\nextframe
\end{fr}

\begin{fr}
\ans{The code. \textbf{A search that succeeded would have refuted a proof}, so
the draws cannot test the theorem \dash{} only whether the program implements
it.}

Stated plainly, because everything in section 5 rests on it:

\medskip
\noindent
\textbf{In a space of $d$ dimensions, at most $d$ vectors can be linearly
independent.} Any $d + 1$ of them are dependent, whatever they are, however
they were produced, and whatever they were trained on.

Notice what that does not need. No inner product, no angles, no notion of
distance \dash{} none of which this program has been given. It follows from
closure and counting alone, which is why it survives every choice of how to
measure the space, and why Program~\ref{prog:P05} can arrive later without
disturbing it.

One word in that sentence is doing damage. \enquote{Dependent} sounds like a
fault. Are the eight vectors of section 2 a defective set?
\nextframe
\end{fr}

\begin{fr}
\ans{No. They are a perfectly good eight vectors that reach two directions.}
A dependent set is not a broken set. The eight vectors of section 2 are a
perfectly good eight vectors; they simply reach two directions rather than
eight. Dependence says something about the \emph{list}, not about the members.

And the redundancy is often the point. A list carrying more vectors than its
span has dimensions is exactly what you want when the extra ones are cheap and
the arrangement is what carries the information \dash{} which is the shape of
the argument section 5 will not quite be able to finish.
\end{fr}

% ============================================================
\section{Basis and dimension}
\label{sec:P04-basis}
\index{basis}
\index{vector!dimension}

\begin{fr}
Two properties have been introduced separately, and putting them together gives
the object the whole subject is built on.

A \emph{basis} for a space is a list of vectors that is \textbf{independent}
and \textbf{spans} it: enough vectors to reach everything, and not one more
than necessary.

Both halves are load-bearing. Drop a vector from a basis and it stops spanning;
add one and it stops being independent. A basis is a list at exactly the size
where those two failures meet.

A space has infinitely many bases. Write down two of them for the same space.
Must they contain the same number of vectors?
\nextframe
\end{fr}

\begin{fr}
\ans{Yes \dash{} and it is the fact that makes the word \emph{dimension} mean
anything at all.}

\medskip
\noindent
\textbf{Every basis for a given space has the same number of vectors.}

\medskip
\noindent
Most readers answer yes without hesitating, which is the right answer arrived
at the wrong way: it is not obvious and it is not a definition \dash{} it is a
theorem, and it is the one this section is really about. If two bases could
differ in size, the word \enquote{dimension} would name nothing.

\begin{rigourbox}
\textbf{Not proved here.} The proof is short and it is the exchange argument:
given two bases, swap the members of one into the other one at a time, showing
at each step that the result still spans. It is in any first course in linear
algebra. What this program needs is the statement, and the statement is what
makes the next paragraph a definition rather than an ambiguity.
\end{rigourbox}

Notice what it does not say. It does not say a space has one basis, or that
any basis is preferred \dash{} this program has already shown that a space has
many. It says only that they all have the same size, which is the weakest
statement a definition could be built on and exactly the one a definition
needs.

That theorem is what licenses a definition rather than an ambiguity, so write
the definition down: the \emph{dimension} of a space is what?
\dotline
\nextframe
\end{fr}

\begin{fr}
\ans{The size of any of its bases \dash{} \emph{any}, because the theorem says
they all have the same one.}

So the four words of this program's title are one idea seen four ways.

\begin{center}
\begin{tabularx}{\linewidth}{@{}lX@{}}
\toprule
span & everything the two operations reach \\
independence & no vector is redundant \\
basis & both at once, at the exact size where they meet \\
dimension & that size, which does not depend on which basis \\
\bottomrule
\end{tabularx}
\end{center}

Each is a sentence about linear combinations, and each is the previous one with
one condition added or removed. Learning them as four definitions is harder
than learning them as one.
\end{fr}

\begin{fr}
Now the property that matters most in practice and gets stated least.

A basis is \textbf{not unique}. The standard axes of the plane are a basis; so
is any other pair of independent vectors, of which there are infinitely many.

Take the point $(\val{p04.basis.x}, \val{p04.basis.y})$ and write it against a
pair of axes rotated by $\val{p04.basis.angle}$ degrees. Its coordinates become
$(\val{p04.basis.c1}, \val{p04.basis.c2})$.

Two things have happened and two have not. Say which is which.
\nextframe
\end{fr}

\begin{fr}
\begin{ansblock}
The coordinates changed. The vector, and its length, did not.
\end{ansblock}

$\sqrt{\val{p04.basis.x}^{2} + \val{p04.basis.y}^{2}} = \val{p04.basis.len}$,
and the new pair gives the same $\val{p04.basis.len}$ \dash{} which the script
asserts, because that is what makes it a change of basis rather than a
different vector.

\textbf{Coordinates are a property of the basis. The vector is a property of the
space.} Everything you would want to say about a vector \dash{} its length, its
relation to another vector, whether a set of them is independent \dash{} is in
the second category and survives any change of basis.
\end{fr}

\mermaidfig{p04-basis-not-unique}{The same vector against two sets of axes, and
what the change does not touch.}{a basis is not unique}

\begin{fr}
\begin{aibox}
Which is what \enquote{embedding dimension $\val{p04.embed}$} does and does not
promise.

It \emph{does} say the embeddings live in a space of $\val{p04.embed}$
dimensions, so at most $\val{p04.embed}$ of them can be independent, and any
list of $\val{p04.embed}$ independent ones is a basis for that space.

It says nothing about the coordinates being meaningful individually. A
component is a coordinate against the basis training happened to land on, and
another training run \dash{} or the same one with a different seed \dash{} lands
on a different basis for a space that may be doing the same job. That is why
\enquote{neuron 2317 detects X} is a claim about a basis and not about a model,
and why comparing two models component by component is comparing two arbitrary
choices of axes.
\end{aibox}
\end{fr}

\begin{fr}
One consequence worth stating before the last section, because it is the
positive form of the same fact.

If coordinates are basis-dependent and the space is not, then \emph{a method
that works in any basis is measuring something real, and a method that needs a
particular basis is measuring your choice of axes}.

That single test disposes of a surprising amount. Comparing two vectors by
length or by whether they are independent: basis-free. Reading meaning off one
component: not. Rotating a model's representations and asking whether it still
works: a test of exactly this distinction. The rotation is itself a matrix,
which is Program~\ref{prog:P06}'s subject; this book does not run the test.
\end{fr}

% ============================================================
\section{What the counting forces, and what it does not}
\label{sec:P04-counting}
\index{vector!dimension}
\index{counting!embedding bound}

\begin{fr}
The payoff, and it is arithmetic you can do in your head once the vocabulary is
in place.

A model has a vocabulary of $\val{p04.vocab}$ tokens and embeds each one as a
vector of $\val{p04.embed}$ components. The rows of that table are
$\val{p04.vocab}$ vectors in a space of $\val{p04.embed}$ dimensions.

At least how many of them are linear combinations of the others?
\nextframe
\end{fr}

\begin{fr}
\ans{$\val{p04.dependent}$, which is $\val{p04.dependent.pct}\%$ of them.}
$\val{p04.vocab} - \val{p04.embed} = \val{p04.dependent}$, and there is nothing
else to it. At most $\val{p04.embed}$ can be independent; the rest are
combinations of some independent subset.

At the $\val{f09.dim}$ dimensions Program~\ref{prog:F09} reasoned about, the
same vocabulary leaves $\val{p04.f09.dependent}$ dependent.

\begin{warning}
\textbf{No amount of training changes this and no architecture avoids it.} It
is not a statement about how well the embeddings were learnt, or about how much
data was used, or about the loss. It follows from two integers, and the only
way to move it is to change one of them.

That is worth dwelling on because it is unusual. Most claims about a model are
empirical and provisional. This one is neither.
\end{warning}
\end{fr}

\begin{fr}
So here is the question it forces, and it is where this program has to be
careful.

The model plainly distinguishes far more than $\val{p04.embed}$ things. It
tells $\val{p04.vocab}$ tokens apart, and beyond tokens it appears to represent
concepts, styles, syntactic roles and much else, in numbers that nobody has
counted but which are not small.

Does the counting bound say it cannot?
\nextframe
\end{fr}

\begin{fr}
\begin{ansblock}
No. It says it cannot give each of them its own \emph{independent direction}.
\end{ansblock}

Those are different claims and the gap between them is the whole of the
interesting question.

A space of $\val{p04.embed}$ dimensions holds exactly $\val{p04.embed}$
mutually independent directions and not one more \dash{} that is the theorem.
It holds vastly more directions that are \emph{almost} independent, and how
many is a question this program cannot ask, because asking it needs a way to
measure \enquote{almost}, which means angles, which is
Program~\ref{prog:P05}'s.

\textbf{Relaxing \enquote{exactly} to \enquote{nearly} is not a technicality
here. It is the entire subject.}
\end{fr}

\begin{fr}
Two accounts of what a model does with that gap are widely discussed, and this
program states both as what they are.

The \emph{linear representation} account holds that human-interpretable
features correspond to directions in the representation space, so that a
feature can be added, removed or measured by moving along a vector. The
\emph{superposition} account holds that a network represents more features than
it has dimensions by packing them into directions that are nearly but not
exactly independent, tolerating the interference that causes because most
features are rare and rarely co-occur.

\begin{warning}
\textbf{Both are hypotheses about structure, and this book states them as
hypotheses.} They are motivated, they have evidence behind them, and neither is
a theorem. What \emph{is} a theorem is the counting bound above, which is why
the two are so often confused: the bound is used as though it established the
account, and it does not. It establishes only that some such account is needed.
\end{warning}

A hypothesis that predicts nothing is a slogan. What would each of the two have
to predict for a test to tell them apart?
\nextframe
\end{fr}

\begin{fr}
\begin{ansblock}
Something that could come out the other way.
\end{ansblock}

The linear account predicts that a feature's presence can be changed by adding
a fixed vector, and that the same vector works across inputs. That is testable
and it is tested. The superposition account predicts a specific kind of
interference \dash{} that features which never co-occur can share
near-directions cheaply and features which do co-occur cannot \dash{} and it
predicts that widening the space should reduce the packing.

\begin{rigourbox}
\textbf{This book does not run either test.} What it can do is keep the
categories apart: a counting bound that is proved, an account of structure that
is proposed, and a measurement that would discriminate between accounts and has
not been made here. Program~\ref{prog:P05} supplies the missing piece of
vocabulary \dash{} what \enquote{nearly independent} means as a number \dash{}
and measures how much of it high dimension gives you for free.
\end{rigourbox}
\end{fr}

\begin{fr}
The last frame, and it is about what has just been bought.

Four words, one idea, and a bound that follows from counting alone. Nothing in
this program needed a length, an angle or a distance, and that is not an
accident: it is why the bound holds under every way of measuring the space you
might later choose.

Program~\ref{prog:F09} said its right-hand column \dash{} how vectors are
arranged, what a typical pair looks like, how much of the space is near the
middle \dash{} was a different class of question and deferred it.
Program~\ref{prog:P05} answers it, and it starts by giving the space the one
thing this program deliberately withheld: a way to multiply two vectors and get
a number back.
\end{fr}

% ============================================================
\begin{summarybox}
\sumitem{1--3}{\result{\textbf{A vector space has two operations and one condition that
  ever fails}: add any two members or scale any member, and the answer must
  land back in the set. The vectors in $\R^{3}$ with first component $0$ are a
  space; those with first component $1$ are not, and the quickest test is that
  \textbf{every space contains the zero vector}}.}
\sumitem{4--5}{\result{The embeddings a model \emph{produces} are a subset of a space,
  not a space \dash{} which is why \enquote{there is a direction for X} is
  ambiguous: the direction exists by construction, and whether the model uses
  it is the question. And \textbf{\enquote{vector} is a role, not a shape}:
  polynomials and functions are vectors, so everything proved from closure
  applies to them unchanged}.}
\sumitem{6--8}{\result{A \textbf{subspace} is a subset that is a space in its own right.
  Closure must hold for \emph{every} pair: the union of the three axes is
  closed under scaling and fails under addition, which is the shape of nearly
  every failed argument. Counting members cannot say which subspace is bigger,
  because both hold infinitely many}.}
\sumitem{9--11}{\result{The \textbf{span} is everything the two operations reach from a
  starting handful, and it is always a subspace. Two vectors in $\R^{3}$ span a
  plane through the origin \dash{} \emph{unless} one is a multiple of the other,
  which is dependence arriving before it is named. \textbf{Two vectors do not
  always give two directions.}}}
\sumitem{12--14}{\result{\textbf{Three numbers, and only the last is the dimension}:
  how many vectors are in the list ($\val{p04.family.count}$), how long each is
  ($\val{p04.family.ambient}$), and how many independent directions they reach
  ($\val{p04.family.dim}$). The first two describe how the data was written
  down; the third describes the data}.}
\sumitem{15}{\result{A list is \textbf{dependent} when one member is a combination
  of the others, equivalently when the zero vector has a non-trivial
  combination; \textbf{independent} when neither. The two forms are the same
  statement divided through by a coefficient}.}
\sumitem{16--18}{\result{\textbf{Rank saturates.} Random vectors in
  $\val{p04.sat.dim}$ dimensions reach $\val{p04.sat.rank.2}$,
  $\val{p04.sat.rank.4}$, $\val{p04.sat.rank.8}$ directions \dash{} and then
  $\val{p04.sat.rank.12}$, $\val{p04.sat.rank.20}$ and
  $\val{p04.sat.rank.40}$, because each new vector either adds a dimension or
  adds nothing, and the space has only so many to give. The
  $\val{p04.sat.trials}$ random draws are testing the code: a counterexample
  would refute a proof}.}
\sumitem{19--20}{\result{\textbf{In $d$ dimensions at most $d$ vectors are
  independent}, whatever they are and whatever they were trained on \dash{} and
  it needs no inner product, no angle and no distance, which is why it survives
  every later choice of how to measure. \emph{Dependent} is not a defect: it
  says something about the list, not the members}.}
\sumitem{21--23}{\result{A \textbf{basis} is independent and spanning \dash{} the exact
  size at which dropping one stops it spanning and adding one stops it being
  independent. \textbf{Every basis for a space has the same size}, which is a
  theorem and is what lets \emph{dimension} name anything. Span, independence,
  basis and dimension are one idea seen four ways}.}
\sumitem{24--25}{\result{\textbf{A basis is not unique.}
  $(\val{p04.basis.x}, \val{p04.basis.y})$ against axes turned by
  $\val{p04.basis.angle}$ degrees becomes
  $(\val{p04.basis.c1}, \val{p04.basis.c2})$, and its length is
  $\val{p04.basis.len}$ either way. \textbf{Coordinates belong to the basis;
  the vector belongs to the space.}}}
\sumitem{26--27}{\result{So \enquote{embedding dimension $\val{p04.embed}$} promises a
  space and a bound, and promises nothing about a component meaning anything on
  its own \dash{} another run lands on another basis. The working test:
  \textbf{a method that works in any basis measures something real; one that
  needs a particular basis measures your choice of axes.}}}
\sumitem{28--29}{\result{\textbf{The counting fact.} $\val{p04.vocab}$ tokens in
  $\val{p04.embed}$ dimensions leaves at least $\val{p04.dependent}$ of the
  embeddings \dash{} $\val{p04.dependent.pct}\%$ \dash{} as combinations of the
  others; at Program~\ref{prog:F09}'s $\val{f09.dim}$ dimensions it is
  $\val{p04.f09.dependent}$. \textbf{No training changes it and no architecture
  avoids it}, which makes it unusual: most claims about a model are provisional
  and this one is not}.}
\sumitem{30--31}{\result{It does \emph{not} say the model cannot distinguish that many
  things \dash{} only that it cannot give each its own \textbf{independent}
  direction. A space of $\val{p04.embed}$ dimensions holds exactly that many
  mutually independent directions and vastly more \emph{nearly} independent
  ones, and \textbf{the relaxation from exactly to nearly is the entire
  subject}}.}
\sumitem{32--34}{\result{The \textbf{linear representation} and \textbf{superposition}
  accounts are \emph{hypotheses about structure}, stated as hypotheses: the
  counting bound is a theorem and is often used as though it established them,
  when it establishes only that some such account is needed. What would
  discriminate between them is a measurement of \enquote{nearly}, which needs
  an angle \dash{} the one thing this program deliberately withheld, and
  Program~\ref{prog:P05}'s opening move}.}
\end{summarybox}

\canyou

% ============================================================
\begin{testexercises}
\item State the two operations of a vector space and the condition their
  results must satisfy.
  \answerto{Addition of two members and multiplication of a member by a scalar;
  both results must lie in the set. That is closure, and it is the condition
  that fails in practice.}
\item Say whether the set of vectors in $\R^{4}$ whose components sum to zero
  is a subspace, with the argument.
  \answerto{Yes. Add two such vectors and the sums add, giving zero; scale one
  and the sum scales, giving zero. It contains the zero vector, as it must.}
\item Give a set of vectors closed under scaling but not under addition.
  \answerto{The union of the coordinate axes: any multiple of $(a, 0, 0)$ still
  has one non-zero component, but $(1, 0, 0) + (0, 1, 0)$ has two.}
\item $\val{p04.family.count}$ vectors of $\val{p04.family.ambient}$ components
  each span $\val{p04.family.dim}$ dimensions. Say which of the three numbers
  is the dimension and what the other two are.
  \answerto{$\val{p04.family.dim}$. The other two are how many vectors were
  written down and how long each one is \dash{} facts about the notation rather
  than about the span.}
\item Say why a list of $d + 1$ vectors in $d$ dimensions cannot be
  independent, without computing anything.
  \answerto{Each vector either adds a dimension to the span or adds nothing.
  The span sits inside a $d$-dimensional space, so it can grow at most $d$
  times, and the $(d+1)$-th vector has nothing left to add.}
\item Define a basis, and say what goes wrong if you remove a vector from one
  and if you add one to it.
  \answerto{Independent and spanning. Remove one and it no longer spans;
  add one and it is no longer independent. A basis sits exactly where those two
  failures meet.}
\item A vector has these coordinates against two different bases:
  \[ (\val{p04.basis.x}, \val{p04.basis.y})
     \qquad\text{and}\qquad
     (\val{p04.basis.c1}, \val{p04.basis.c2}) \]
  Say what this tells you about reading meaning off a single component.
  \answerto{That it is a statement about the basis rather than about the
  vector. Both coordinate pairs describe the same vector of length
  $\val{p04.basis.len}$, and nothing distinguishes one pair as the real one.}
\item A model embeds $\val{p04.vocab}$ tokens in $\val{p04.embed}$ dimensions.
  Give the number that must be dependent and say what would change it.
  \answerto{At least $\val{p04.dependent}$. Only changing one of the two
  integers \dash{} a smaller vocabulary or a wider embedding. No amount of
  training does.}
\item Say what the counting bound proves about a model's ability to represent
  many concepts, and what it does not.
  \answerto{It proves the model cannot give each concept its own independent
  direction. It does not prove the model cannot distinguish them, because
  nearly-independent directions are far more numerous and the bound says
  nothing about those.}
\item Say why the linear representation and superposition accounts are
  described in this program as hypotheses.
  \answerto{Because they are claims about how a trained network is structured,
  supported by evidence and not established by proof. The counting bound is the
  theorem, and it shows only that some account of this kind is needed.}
\end{testexercises}

% ============================================================
\begin{furtherproblems}
\item Show that the span of any set of vectors is a subspace, directly from the
  definition.
  \answerto{A sum of two linear combinations of the $v_i$ collects into one
  linear combination of the $v_i$, and a scalar multiple of one does too. Both
  results are in the span, so it is closed under both operations.}
\item Show that the two definitions of independence are equivalent.
  \answerto{If $a_1 v_1 + \dots + a_k v_k = 0$ with some $a_j \ne 0$, divide by
  $a_j$ and rearrange to write $v_j$ as a combination of the rest. Conversely,
  if $v_j$ is such a combination, moving it to the other side gives a
  non-trivial combination equal to zero.}
\item A list of vectors contains the zero vector. Say whether it can be
  independent.
  \answerto{No. Take the coefficient of the zero vector to be $1$ and every
  other coefficient to be $0$: that is a non-trivial combination equal to zero.
  A list containing the zero vector is dependent whatever else is in it.}
\item Two subspaces of $\R^{5}$ have dimensions $3$ and $4$. Say what you can
  conclude about their intersection.
  \answerto{It has dimension at least $2$. Their dimensions sum to $7$ in a
  space of $5$, so they overlap by at least the excess \dash{} in particular
  they cannot meet only at the origin, and there is a whole plane in both.}
\item Explain why the counting bound needs no inner product, and say what that
  buys.
  \answerto{Independence is defined by linear combinations alone, and the
  saturation argument counts dimensions rather than measuring anything. So the
  bound holds under every possible way of equipping the space with lengths and
  angles \dash{} it cannot be escaped by choosing a different metric.}
\item A model's representations are multiplied throughout by a fixed rotation.
  Say which of these change: a component's value, a vector's length, whether a
  set is independent, the dimension of the span.
  \answerto{Only the component's value. The other three are basis-free, which
  is exactly the test the program gives for whether a claim is about the model
  or about the axes.}
\item Suppose a model's $\val{p04.vocab}$ embeddings were made exactly
  independent by widening the space. Give the width required, and say what it
  would cost using Program~\ref{prog:P03}'s arithmetic.
  \answerto{$\val{p04.vocab}$ dimensions. The embedding table alone would then
  have $\val{p04.vocab}^{2}$ parameters \dash{} four hundred million \dash{}
  and every layer's width would follow it, which is why nobody does this and
  why the packing question exists at all.}
\item Argue that the counting bound cannot by itself distinguish between the
  superposition account and an account in which the model simply represents
  fewer features than people think.
  \answerto{The bound constrains independent directions and says nothing about
  how many features there are. Both accounts are consistent with it: one posits
  many features packed into near-directions, the other posits few features in
  independent ones. Discriminating between them needs a measurement of the
  arrangement, not a count.}
\end{furtherproblems}
